% Part: many-valued-logic
% Chapter: three-valued-logics
% Section: kleene
%READERNOTE{SOL6-A141: दोन्ही क्लिनी मॅट्रिक्स येथे मानक L0 च्या असत्यता-स्थिर नसलेल्या चार-संयोजकीय खंडावर आहेत. सर्व चलांना U देणाऱ्या मूल्यनिर्णयनावरील सर्वतः सत्य सूत्र नसण्याचा पुरावा या व्याप्तीत लागू होतो.}

\documentclass[../../../include/open-logic-section]{subfiles}

\begin{document}

\olfileid{mvl}{thr}{skl}

\olsection{क्लिनी तर्कशास्त्रे}

Stephen Kleene यांनी संगणन प्रक्रियांच्या परिणामांच्या दृष्टीने
सत्यतामूल्यांचा विचार करणाऱ्या तर्कशास्त्रातून प्रेरित होऊन दोन
त्रिमूल्य तर्कशास्त्रे मांडली: एखादी प्रक्रिया $\True$ किंवा
$\False$ हा परिणाम देऊ शकते; पण ती समाप्तही होणार नाही असे
घडू शकते. त्या वेळी संबंधित सत्यतामूल्य अपरिभाषित असते;
ते सत्यतामूल्य~$\Undef$ ने दर्शवले जाते.

विधान~$!A$ चे निषेधन संगणित करण्यासाठी आधी~$!A$ चे मूल्य
काढायचे आणि मग त्याच्या विरुद्ध मूल्याचा परिणाम द्यायचा.
$!A$ चे संगणन समाप्त होत नसेल, तर ही संपूर्ण प्रक्रियाही
समाप्त होत नाही; म्हणून $\Undef$ चे निषेधन $\Undef$ आहे.

संधी $!A \land !B$ संगणित करण्याचे दोन मार्ग आहेत: आधी~$!A$
आणि नंतर $!B$ संगणित करता येईल. दोन्हींचा परिणाम~$\True$
असेल तर संधीचा परिणाम $\True$, अन्यथा $\False$ असेल.
दोन्हींपैकी कोणतेही संगणन थांबले नाही, तर संपूर्ण प्रक्रियाही
थांबत नाही. त्यामुळे या पद्धतीत संधीचा एक घटक अपरिभाषित
असेल, तर संधीही अपरिभाषित असते. विसंधीबाबतही हेच लागू होते.

पण $!A$ आणि $!B$ यांचे मूल्यांकन समांतरपणे करता आले, तर
अधिक चांगले करता येते. दोन प्रक्रियांपैकी एक थांबून $\False$
हा परिणाम देताच आपण थांबू शकतो, कारण संधी असत्य असलीच
पाहिजे. म्हणून या पद्धतीत एक घटक असत्य असेल, तर दुसरा
अपरिभाषित असला तरी संधी असत्य असते. त्याचप्रमाणे विसंधीचे
समांतर संगणन करताना दोन घटकांपैकी एक सत्य परिणाम देताच
थांबता येते; विसंधी सत्य असलीच पाहिजे. त्यामुळे संयुक्त
विधानाचा एक घटक अपरिभाषित असला तरी त्याचा परिणाम कळू
शकतो. या अर्थानुसार $\Undef$ ला ``अपरिभाषित'' याऐवजी
``अज्ञात'' असे वाचता येईल.

या दोन अर्थांमधून क्लिनी यांची प्रबळ आणि दुर्बल तर्कशास्त्रे
मिळतात. सशर्त विधानाची व्याख्या $\lnot !A \lor !B$ शी
समतुल्य अशी केली जाते.

\begin{defn}
\emph{प्रबळ क्लिनी तर्कशास्त्र}~$\LogKs$ पुढील मॅट्रिक्सने
परिभाषित होते:
\begin{enumerate}
  \item मानक विधानीय भाषा $\Lang L_0$ चा केवळ
  $\lnot$, $\land$, $\lor$, $\lif$ हे संयोजक असलेला खंड.
  \item सत्यतामूल्यांचा संच $V = \{\True, \Undef, \False\}$.
  \item $\True$ हेच एकमेव निर्दिष्ट मूल्य आहे, म्हणजे $V^+ = \{\True\}$.
  \item सत्यता-फलने पुढील कोष्टकांनी दिली आहेत:
  \begin{center}
    \begin{tabular}{c|c}
      $\tf{\lnot}$ & \\
      \hline
      $\True$ & $\False$ \\
      $\Undef$ & $\Undef$ \\
      $\False$ & $\True$
    \end{tabular}
    \quad
    \begin{tabular}{c|ccc}
      $\tf{\land}[\LogKs]$ & $\True$ & $\Undef$ & $\False$ \\
      \hline
      $\True$ & $\True$ & $\Undef$ & $\False$ \\
      $\Undef$ & $\Undef$ & $\Undef$ & $\False$\\
      $\False$ & $\False$ & $\False$ & $\False$
    \end{tabular}
    \\[2ex]
    \begin{tabular}{c|ccc}
      $\tf{\lor}[\LogKs]$ & $\True$ & $\Undef$ & $\False$ \\
      \hline
      $\True$ & $\True$ & $\True$ & $\True$ \\
      $\Undef$ & $\True$ & $\Undef$ & $\Undef$ \\
      $\False$ & $\True$ & $\Undef$ & $\False$
    \end{tabular}
    \quad
    \begin{tabular}{c|ccc}
      $\tf{\lif}[\LogKs]$ & $\True$ & $\Undef$ & $\False$ \\
      \hline
      $\True$ & $\True$ & $\Undef$ & $\False$ \\
      $\Undef$ & $\True$ & $\Undef$ & $\Undef$  \\
      $\False$ & $\True$ & $\True$ & $\True$
    \end{tabular}
  \end{center}
\end{enumerate}
\end{defn}

\begin{defn}
\emph{दुर्बल क्लिनी तर्कशास्त्र}~$\LogKw$ पुढील मॅट्रिक्सने
परिभाषित होते:
\begin{enumerate}
  \item मानक विधानीय भाषा $\Lang L_0$ चा केवळ
  $\lnot$, $\land$, $\lor$, $\lif$ हे संयोजक असलेला खंड.
  \item सत्यतामूल्यांचा संच $V = \{\True, \Undef, \False\}$.
  \item $\True$ हेच एकमेव निर्दिष्ट मूल्य आहे, म्हणजे $V^+ = \{\True\}$.
  \item सत्यता-फलने पुढील कोष्टकांनी दिली आहेत:
  \begin{center}
    \begin{tabular}{c|c}
      $\tf{\lnot}$ & \\
      \hline
      $\True$ & $\False$ \\
      $\Undef$ & $\Undef$ \\
      $\False$ & $\True$
    \end{tabular}
    \quad
    \begin{tabular}{c|ccc}
      $\tf{\land}[\LogKw]$ & $\True$ & $\Undef$ & $\False$ \\
      \hline
      $\True$ & $\True$ & $\Undef$ & $\False$ \\
      $\Undef$ & $\Undef$ & $\Undef$ & $\Undef$\\
      $\False$ & $\False$ & $\Undef$ & $\False$
    \end{tabular}
    \\[2ex]
    \begin{tabular}{c|ccc}
      $\tf{\lor}[\LogKw]$ & $\True$ & $\Undef$ & $\False$ \\
      \hline
      $\True$ & $\True$ & $\Undef$ & $\True$ \\
      $\Undef$ & $\Undef$ & $\Undef$ & $\Undef$ \\
      $\False$ & $\True$ & $\Undef$ & $\False$
    \end{tabular}
    \quad
    \begin{tabular}{c|ccc}
      $\tf{\lif}[\LogKw]$ & $\True$ & $\Undef$ & $\False$ \\
      \hline
      $\True$ & $\True$ & $\Undef$ & $\False$ \\
      $\Undef$ & $\Undef$ & $\Undef$ & $\Undef$  \\
      $\False$ & $\True$ & $\Undef$ & $\True$
    \end{tabular}
  \end{center}
\end{enumerate}
\end{defn}

\begin{prop}
  $\LogKs$ आणि $\LogKw$ यांत कोणतेही सर्वतः सत्य सूत्र नाही.
\end{prop}

\begin{proof}
  सर्व !!{propositional variable}s~$p$ साठी
  $\pAssign v(p) = \Undef$ असेल, तर कोणत्याही
  सूत्र~$!A$ चे सत्यतामूल्य~$\pValue v(!A) =
  \Undef$ असेल; कारण दोन्ही तर्कशास्त्रांत
  \[
    \tf{\lnot}(\Undef) = \tf{\lor}(\Undef, \Undef) = \tf{\land}(\Undef,
  \Undef) = \tf{\lif}(\Undef, \Undef) = \Undef
  \]
  असे आहे. $\LogKs$ आणि $\LogKw$ या दोन्हींमध्ये
  $\Undef \notin V^+$ असल्यामुळे या !!{valuation} मध्ये
  $!A$ चे सत्यतामूल्य निर्दिष्ट नसेल.
\end{proof}

दुर्बल आणि प्रबळ क्लिनी तर्कशास्त्रांत सर्वतः सत्य सूत्रे
नसली, तरी त्यांचे तार्किक निष्पन्नता संबंध अतुच्छ आहेत.

\begin{prob}
  पुढीलपैकी कोणते संबंध (a) प्रबळ आणि (b) दुर्बल क्लिनी
  तर्कशास्त्रात खरे ठरतात? प्रत्येकासाठी सत्यताकोष्टक द्या.
  \begin{enumerate}
    \item $p, p \lif q \Entails q$
    \item $p \lor q, \lnot p \Entails q$
    \item $p \land q \Entails p$
    \item $p \Entails p \land p$
    \item $p \Entails p \lor q$
  \end{enumerate}
\end{prob}

Dmitry Bochvar यांनी $\Undef$ चा अर्थ ``अर्थहीन'' असा घेतला
आणि विरोधाभास निर्माण करणाऱ्या विधानांना~$\Undef$ हे मूल्य
देऊन खोटारड्याच्या विरोधाभासासारखे विरोधाभास सोडवण्याचा
प्रयत्न केला. त्यांनी मूलतः दुर्बल क्लिनी तर्कशास्त्राला
अतिरिक्त संयोजके जोडून तयार होणारे तर्कशास्त्र मांडले.
त्यांतील दोन संयोजके ``बाह्य निषेधन'' आणि ``अपरिभाषित
आहे'' हा परिकर्मी अशी आहेत:
\begin{center}
  \begin{tabular}{c|c}
    $\tf{\sim}$ & \\
    \hline
    $\True$ & $\False$ \\
    $\Undef$ & $\True$ \\
    $\False$ & $\True$
  \end{tabular}
\quad
  \begin{tabular}{c|c}
    $\tf{+}$ & \\
    \hline
    $\True$ & $\False$ \\
    $\Undef$ & $\True$ \\
    $\False$ & $\False$
  \end{tabular}
\end{center}

\begin{prob}
  Bochvar यांच्या तर्कशास्त्रात $\lnot$ आणि~$+$ यांच्या
  साहाय्याने $\sim$ ची व्याख्या करता येते का? म्हणजे, फक्त
  !!{propositional
  variable}~$p$ असलेले, $\sim$ नसलेले
  आणि नेहमी~$\mathord{\sim}p$ इतकेच सत्यतामूल्य घेणारे
  !!a{formula} शोधा. तुमचे उत्तर योग्य आहे हे
  सत्यताकोष्टकाने दाखवा.
\end{prob}


\end{document}
